\documentclass[11pt]{article}
\usepackage{docstyle}
\renewcommand\sheetname{In-class work \textperiodcentered{} Human Reproduction}
\begin{document}
\sheethead{In-class work + Mini mock}{Human Reproduction \textperiodcentered{} Year 9 Science}{Each question: A / M / E}
Part A after Loop 1, Part B after Loop 2, Part C after Loop 3, Part D after Loop 4. Part E is the mini mock at the end of the lesson.


\section{Part A \textperiodcentered{} Sex cells and the two systems}
\Q The diagram shows the female reproductive system. Write the letter of the part for each sentence.\lvl{A}
\par\vspace{1mm}
\begin{minipage}[t]{0.50\linewidth}\vspace{0pt}\FigFemaleFront[5mm]\end{minipage}\hfill
\begin{minipage}[t]{0.46\linewidth}\vspace{4mm}\raggedright
(a) Eggs are made here. \hfill\blank[14mm]{A}\par\vspace{3.5mm}
(b) Fertilisation happens here. \hfill\blank[14mm]{B}\par\vspace{3.5mm}
(c) An embryo implants in this. \hfill\blank[14mm]{D}\par\vspace{3.5mm}
(d) Sperm are received here. \hfill\blank[14mm]{F}\par\vspace{3.5mm}
(e) This is a ring of muscle. \hfill\blank[14mm]{E}
\end{minipage}\par\vspace{2mm}

\Q Complete the sentences about the male reproductive system.\lvl{A}
\par Sperm are made in the \blank[26mm]{testes}, which are held outside the body in the \blank[26mm]{scrotum}.
The sperm are stored in the \blank[28mm]{epididymis}. They travel along the sperm \blank[20mm]{duct}, where glands add fluid to make \blank[22mm]{semen}.
The semen leaves the body through the \blank[26mm]{urethra}.\par\vspace{2mm}

\Q The middle piece of a sperm cell is packed with mitochondria. Explain how this helps the sperm cell to do its job.\lvl{M}
\alines{3}{Mitochondria release energy (by respiration). The tail uses this energy to swim, so the sperm can travel through the uterus to the egg in the oviduct.}

\Q An egg cell is one of the largest cells in the body. A sperm cell is one of the smallest. Explain why each size suits the job of the cell.\lvl{E}
\alines{4}{The egg is large because its cytoplasm holds a food store, which feeds the embryo for its first cell divisions, before it implants. The sperm is small and streamlined, so it needs little energy to swim and millions of them can be made.}

\section{Part B \textperiodcentered{} Hormones and the menstrual cycle}
\Q Name the four stages of the menstrual cycle, in order, starting with the period.\lvl{A}
\par 1. \blank[32mm]{menstrual} \quad 2. \blank[32mm]{follicular} \quad 3. \blank[32mm]{ovulation} \quad 4. \blank[30mm]{luteal}\par\vspace{2mm}

\Q Write the name of the hormone that does each job.\lvl{A}
\par (a) It makes an egg mature in the ovary. \hfill\blank[44mm]{FSH}
\par (b) It triggers ovulation. \hfill\blank[44mm]{LH}
\par (c) It builds up the lining of the uterus after a period. \hfill\blank[44mm]{oestrogen}
\par (d) It keeps the lining of the uterus thick after ovulation. \hfill\blank[44mm]{progesterone}\par\vspace{2mm}

\Q The table shows the level of progesterone in a woman's blood during one cycle.
\par\vspace{1mm}
\begin{tabular}{@{}l *{8}{c}@{}}
\toprule
Day of the cycle & 2 & 6 & 10 & 14 & 18 & 22 & 26 & 28\\
Progesterone level (units) & 1 & 1 & 1 & 2 & 9 & 14 & 6 & 2\\
\bottomrule
\end{tabular}\par\vspace{1mm}
\sub{a} Describe how the level of progesterone changes during the cycle.\lvl{M}
\alines{3}{It stays low, at 1 unit, until day 10. After day 14 it rises quickly to a peak of 14 units on day 22. Then it falls to 2 units by day 28.}
\sub{b} Explain what the fall after day 22 tells you about this cycle.\lvl{E}
\alines{3}{The egg was not fertilised, so no embryo implanted. Progesterone falls, so the lining of the uterus is no longer kept thick: it breaks down and a period starts.}

\section{Part C \textperiodcentered{} Fertilisation, chromosomes and twins}
\Q Write these in the order in which they form, earliest first: \textbf{embryo, foetus, sex cells, zygote}.\lvl{A}
\par \blank[34mm]{sex cells} $\rightarrow$ \blank[34mm]{zygote} $\rightarrow$ \blank[34mm]{embryo} $\rightarrow$ \blank[34mm]{foetus}
\par\vspace{2.5mm} Number of chromosomes in a human sperm cell: \blank[16mm]{23} \quad in a human zygote: \blank[16mm]{46}\par\vspace{2mm}

\Q A sperm cell carrying a Y chromosome fertilises an egg cell.\lvl{A / M}
\par\sub{a} Write the two sex chromosomes of the zygote. \hfill\blank[30mm]{XY}
\par\vspace{1mm}\sub{b} State the sex of the baby. \hfill\blank[30mm]{male (a boy)}
\par\vspace{1mm}\sub{c} Explain why it is the sperm cell, not the egg cell, that decides the sex of a baby.
\alines{3}{Every egg cell carries an X chromosome. A sperm cell carries either X or Y. So the only difference comes from the sperm: X gives XX (a girl), Y gives XY (a boy).}

\Q Ana and Bea are twin sisters, but they do not look alike. Explain how these twins formed and why they are not identical.\lvl{E}
\alines{4}{Two eggs were released, and each was fertilised by a different sperm, so there were two zygotes. Each zygote had a different mixture of genes, so the sisters are non-identical. (They are the same sex only by chance.)}

\section{Part D \textperiodcentered{} Pregnancy and birth}
\Q The diagram shows a foetus in the uterus. Write the letter of the part for each sentence.\lvl{A}
\par\vspace{1mm}
\begin{minipage}[t]{0.50\linewidth}\vspace{0pt}\FigWomb[3.9mm]\end{minipage}\hfill
\begin{minipage}[t]{0.46\linewidth}\vspace{5mm}\raggedright
(a) It cushions the foetus. \hfill\blank[14mm]{C}\par\vspace{3.5mm}
(b) It joins the foetus to the placenta. \hfill\blank[14mm]{B}\par\vspace{3.5mm}
(c) Substances pass to and from the mother's blood here. \hfill\blank[14mm]{A}\par\vspace{3.5mm}
(d) It opens (dilates) during labour. \hfill\blank[14mm]{F}
\end{minipage}\par\vspace{2mm}

\Q Urea is a waste product made by the cells of the foetus. Explain how urea leaves the blood of the foetus.\lvl{M}
\alines{3}{The concentration of urea is higher in the foetus's blood than in the mother's blood. In the placenta, urea diffuses across the thin wall of the villi into the mother's blood, which carries it away.}

\Q The table shows the mass of a foetus during pregnancy.
\par\vspace{1mm}
\begin{tabular}{@{}l *{5}{c}@{}}
\toprule
Week of pregnancy & 8 & 16 & 24 & 32 & 40\\
Mass of foetus (g) & 1 & 120 & 640 & 1760 & 3400\\
\bottomrule
\end{tabular}\par\vspace{1mm}
\sub{a} Calculate the increase in mass between week 32 and week 40.\lvl{A}
\alines{1}{$3400 - 1760 = 1640$ g}
\sub{b} Describe how the mass of the foetus changes during pregnancy.\lvl{M}
\alines{3}{The mass increases all the time. It increases slowly at first, to 120 g by week 16, and then faster and faster: from 1760 g at week 32 to 3400 g at week 40.}
\sub{c} A baby born at week 32 usually needs special care in hospital. Use the table to suggest why.\lvl{E}
\alines{3}{The baby is premature. Its mass is only about half the mass of a baby at 40 weeks (1760 g, not 3400 g), and its organs, such as the lungs, have not finished developing.}

\newpage
\section{Part E \textperiodcentered{} Mini mock (12 minutes)}
\Q State what is meant by \textbf{fertilisation}.\lvl{A}
\alines{2}{Fertilisation is when the nucleus of a sperm cell fuses (joins) with the nucleus of an egg cell.}

\Q State the number of chromosomes in:\lvl{A}
\par (a) a human egg cell \blank[18mm]{23} \qquad (b) a human skin cell \blank[18mm]{46}\par\vspace{2mm}

\Q The diagram shows the male reproductive system. Write the letter of the part for each sentence.\lvl{A}
\par\vspace{1mm}
\begin{minipage}[t]{0.44\linewidth}\vspace{0pt}\FigMaleFront[46mm]\end{minipage}\hfill
\begin{minipage}[t]{0.52\linewidth}\vspace{8mm}\raggedright
(a) Sperm are made here. \hfill\blank[14mm]{G}\par\vspace{4mm}
(b) Sperm are stored here. \hfill\blank[14mm]{F}\par\vspace{4mm}
(c) This tube carries urine and semen. \hfill\blank[14mm]{E}\par\vspace{4mm}
(d) This keeps the testes cool. \hfill\blank[14mm]{H}
\end{minipage}\par\vspace{2mm}

\Q Explain why the lining of the uterus becomes thick every month.\lvl{M}
\alines{3}{If an egg is fertilised, the embryo needs somewhere to implant. The thick lining has many blood vessels, which bring the food and oxygen that the embryo needs to grow.}

\Q A student is asked to describe implantation.\par\vspace{1mm}
\samcard{0.9\linewidth}{\flawed{Implantation is when the sperm joins the egg in the uterus. The ball of cells then sinks into the wall of the oviduct.}}\par\vspace{1mm}
\textbf{Sam's answer has 2 errors.} Find each one and write the correct science.\lvl{M}
\alines{4}{1. The sperm joining the egg is fertilisation, and it happens in the oviduct (\Lref{fert}). 2. In implantation the embryo (ball of cells) sinks into the lining of the uterus, not the oviduct (\Lref{stages}).}

\Q A pregnant woman is advised not to drink alcohol. Explain how alcohol in her blood can reach the foetus, and why it is harmful.\lvl{E}
\alines{4}{Alcohol diffuses across the placenta, from the mother's blood (high concentration) into the foetus's blood (low concentration). The umbilical cord carries it to the foetus. The foetus is very small and growing quickly, so even a little alcohol can damage its developing brain.}

\Q The table shows how often IVF treatment led to the birth of a baby, for women of different ages.
\par\vspace{1mm}
\begin{tabular}{@{}l *{5}{c}@{}}
\toprule
Age of woman (years) & under 35 & 35--37 & 38--39 & 40--42 & 43--44\\
Treatments that led to a birth (\%) & 32 & 25 & 19 & 11 & 5\\
\bottomrule
\end{tabular}\par\vspace{1mm}
\sub{a} Describe the pattern in the table.\lvl{M}
\alines{2}{The older the woman, the lower the percentage: it falls from 32\,\% for women under 35 to 5\,\% for women aged 43--44.}
\sub{b} Suggest one reason for this pattern.\lvl{E}
\alines{2}{Older women release fewer ripe, healthy eggs, so fewer eggs are fertilised and fewer embryos implant.}
\end{document}
